% ============================================================
%  PDL--OFN Joint Note N02
%  From Z3 to Three Generations: A Structural Bridge between
%  PDL Leakage Cycles and OFN Fermion Families
%  C. Laubscher, O. I. Evdokimov -- 2026
%
%  STATUS OF THIS FILE: PDL-side draft, circulated to Oleg for
%  the OFN-side contribution (Section 4) and joint revision of
%  Sections 5-6. Do not deposit before both sides are complete
%  and the comparison table (Section 4) has been jointly reviewed.
% ============================================================
\documentclass[11pt,a4paper]{article}

\usepackage[utf8]{inputenc}
\usepackage[T1]{fontenc}
\usepackage[british]{babel}
\usepackage[margin=2.5cm]{geometry}
\usepackage{amsmath,amssymb,amsthm}
\usepackage{booktabs}
\usepackage{array}
\usepackage{graphicx}
\usepackage{float}
\usepackage[font=small,labelfont=bf]{caption}
\usepackage{xcolor}
\usepackage{hyperref}
\usepackage[numbers]{natbib}
\usepackage{enumitem}
\usepackage[activate=false]{microtype}
\usepackage{tcolorbox}
\tcbuselibrary{theorems,skins}

\definecolor{proofblue}{RGB}{0,70,140}
\definecolor{conjecturegreen}{RGB}{0,120,60}
\definecolor{hypothesisorange}{RGB}{200,90,0}
\definecolor{openproblemred}{RGB}{170,0,0}
\definecolor{speculativegrey}{RGB}{90,90,90}

\hypersetup{
  colorlinks=true,
  linkcolor=proofblue,
  citecolor=proofblue,
  urlcolor=proofblue
}

% ---- Theorem environments (PDL house style) ----
\theoremstyle{plain}
\newtheorem{theorem}{Theorem}[section]
\newtheorem{proposition}[theorem]{Proposition}
\newtheorem{corollary}[theorem]{Corollary}
\newtheorem{lemma}[theorem]{Lemma}
\theoremstyle{definition}
\newtheorem{definition}[theorem]{Definition}
\theoremstyle{remark}
\newtheorem{remark}[theorem]{Remark}
\newtheorem{conjecture}[theorem]{Conjecture}
\newtheorem{openproblem}[theorem]{Open Problem}

% ---- tcolorbox environments for cross-framework material ----
\newtcolorbox{pdlbox}[1][]{
  colback=blue!4!white, colframe=proofblue,
  fonttitle=\bfseries, title={PDL (established)}, #1
}
\newtcolorbox{ofnbox}[1][]{
  colback=green!4!white, colframe=conjecturegreen,
  fonttitle=\bfseries, title={OFN (established)}, #1
}
\newtcolorbox{analogybox}[1][]{
  colback=orange!4!white, colframe=hypothesisorange,
  fonttitle=\bfseries, title={Candidate structural analogy --- not a derivation}, #1
}
\newtcolorbox{speculativebox}[1][]{
  colback=gray!4!white, colframe=speculativegrey,
  fonttitle=\bfseries, title={Speculative extension}, #1
}

% ---- Macros ----
\newcommand{\PDL}{\textnormal{PDL}}
\newcommand{\OFN}{\textnormal{OFN}}
\newcommand{\Kfour}{K_4}
\newcommand{\Sfour}{S_4}
\newcommand{\Sthree}{S_3}
\newcommand{\Afour}{A_4}
\newcommand{\Vfour}{V_4}
\newcommand{\Zthree}{\mathbb{Z}_3}
\newcommand{\SUthree}{\mathrm{SU}(3)}
\renewcommand{\Re}{R_e}
\newcommand{\Rtot}{R_{\mathrm{tot}}}
\newcommand{\Rsea}{R_{\mathrm{sea}}}
\newcommand{\rval}{r_{\mathrm{val}}}
\newcommand{\nuu}{n_u}
\newcommand{\ndd}{n_d}
\newcommand{\Dn}{\Delta n}
\newcommand{\Dmiso}{\Delta m_{\mathrm{iso}}}
\newcommand{\kone}{k_1}
\newcommand{\ktwo}{k_2}
\newcommand{\kthree}{k_3}
\newcommand{\pkone}{p_{k_1}}
\newcommand{\Omegatf}{\Omega_{21}}
\newcommand{\Ssr}{S_{\mathrm{sr}}}

% ============================================================
\begin{document}
% ============================================================

\title{\textbf{From \texorpdfstring{$\Zthree$}{Z3} to Three Generations:\\
A Structural Bridge between \PDL\ Leakage Cycles\\
and \OFN\ Fermion Families}\\[6pt]
{\large Joint PDL--OFN Note N02}}

\author{Oleg I.~Evdokimov\thanks{Astronomical Observatory, Kazan State University, Kazan', Tatarstan Republic, Russian Federation. ORCID: \href{https://orcid.org/0009-0005-3624-8504}{0009-0005-3624-8504}.}
\quad\and\quad
C\'edric Laubscher\thanks{Independent researcher, Switzerland. ORCID: \href{https://orcid.org/0009-0004-5415-1098}{0009-0004-5415-1098}. \href{https://cedriclaubscher.ch}{cedriclaubscher.ch}.}}

\date{July 2026 --- \textbf{Draft v0.2, complete joint draft for review}}

\maketitle

% ============================================================
\begin{center}
\fbox{\parbox{0.92\textwidth}{
\centering\textbf{Status of this document}\\[4pt]
\raggedright
This document is a joint draft of N02, integrating the complete \PDL-side (Sections~1--3, 5--6) and the complete \OFN-side (Section~4). The comparison table (Section~3) and open problems (Section~7) have been updated to reflect all results established in July 2026. The document is ready for final joint review before Zenodo deposit.
}}
\end{center}

\vspace{0.5em}

% ============================================================
\begin{abstract}
The joint note N01~\cite{N01} established that the first Betti number $\beta_1(\Kfour) = 3$ (Projective Dynamic Logo, \PDL) and the cyclomatic number $b_1(\Omegatf) = 3$ (Ontological Fundamental Network, \OFN) are the same topological invariant arising independently in two unrelated combinatorial frameworks. Subsequent \PDL\ documents (D57--D59) trace this invariant to a specific algebraic origin internal to \PDL: the quotient $\Afour/\Vfour \cong \Zthree$ in $\mathrm{Sym}(\Kfour) = \Sfour$, which is simultaneously the centre of the colour gauge group $\SUthree$ (D58) and the index set of the three axes of the fundamental representation $\mathbf{3}$ (D59). The present note examines whether this $\Zthree$ structure, established as an unconditional \PDL\ theorem, can be placed in correspondence with the three fermion generations of the \OFN\ framework, thereby providing a partial resolution of the open problem OP-OFN-1. We present the complete \PDL-side derivation chain together with a candidate structural analogy to a numerically verified arithmetic identity linking the \PDL\ hadronic and cosmological sectors ($\nuu - 1 = \pkone = 23$), and we propose a comparison table format, following the working conditions of N01, within which the \OFN-side correspondence --- including the structural self-reference operator $\Ssr$ --- is to be developed jointly.
\end{abstract}

\clearpage

\tableofcontents

\clearpage

% ============================================================
\section{Introduction and motivation}
\label{sec:intro}
% ============================================================

N01~\cite{N01} posed three open problems at the interface of the \PDL\ and \OFN\ frameworks. Of these, OP-OFN-1 --- the formal link between the three \PDL\ leakage cycles and the three \OFN\ fermion generations --- has remained the most resistant, because at the time of N01 neither framework had yet identified the precise algebraic mechanism generating its respective invariant~3.

Since N01, three \PDL\ documents have closed part of this gap from the \PDL\ side. D57~\cite{D57} identified the quotient $\Afour/\Vfour \cong \Zthree$ within $\mathrm{Sym}(\Kfour) = \Sfour$ as a candidate combinatorial origin of the invariant. D58~\cite{D58} proved, as an unconditional algebraic theorem of the \PDL\ axioms C1--C4, that this same quotient $\Zthree$ is exactly the centre of the colour gauge group $\SUthree$. D59~\cite{D59} went further, constructing explicitly the fundamental representation $\mathbf{3}$ of $\SUthree$ on the trace-zero plane $W \subset \mathbb{C}^3$, whose three axes are canonically labelled both by $\Vfour \setminus \{e\}$ (the geometric labelling) and by the triplet $(T_2) = \{\kone, \ktwo, \kthree\}$ of leakage-cycle indices from D51--D52 (the dynamical labelling).

The purpose of the present note is twofold. First, to assemble this \PDL-side chain in a single self-contained account, so that the precise algebraic content of the invariant~3 is available for comparison with its \OFN\ counterpart. Second, to propose --- under the same epistemic discipline as N01 --- a structured format within which the \OFN-side mechanism (in particular the structural self-reference operator $\Ssr$ introduced in the \OFN\ programme~\cite{DTOCOFN}) can be set alongside the \PDL\ result without conflating the two frameworks.

A secondary, more speculative thread is introduced in Section~\ref{sec:speculative}: the arithmetic identity $\nuu - 1 = \pkone = 23$ (D63), which connects the \PDL\ hadronic sector (the proton valence multiplicity $\nuu$) to the \PDL\ cosmological sector (the first leakage-cycle prime $\pkone$). This identity is independent of, and additional to, the main $\Zthree$ chain; it is presented separately and at a lower epistemic register, consistent with its open status (OP-D63-3) in the \PDL\ corpus.

% ============================================================
\section{The PDL side: from \texorpdfstring{$\Afour/\Vfour \cong \Zthree$}{A4/V4 = Z3} to the three axes of \texorpdfstring{$W$}{W}}
\label{sec:pdl_side}
% ============================================================

\subsection{The elementary closure and its automorphisms}

The \PDL\ axioms C1--C4 force the elementary closure to be the complete signed graph $\Kfour$ on four vertices and six edges~\cite{D47}. The automorphism group of $\Kfour$ as an unsigned graph is $\Sfour$, acting on the vertex set by permutation. Two distinguished normal subgroups of $\Sfour$ are central to the chain assembled here: the Klein four-group $\Vfour = \{e, (01)(23), (02)(13), (03)(12)\}$, and the alternating subgroup $\Afour$ of even permutations, with the inclusion $\Vfour \subset \Afour \subset \Sfour$ and indices $[\Afour:\Vfour] = 3$, $[\Sfour:\Afour] = 2$.

\subsection{Three theorems: \texorpdfstring{$\Sthree$}{S3}, \texorpdfstring{$\Zthree$}{Z3}, and the representation \texorpdfstring{$\mathbf{3}$}{3}}

\begin{pdlbox}
\textbf{Theorem (D60).} The physical indistinguishability of the two states of a pulsation pairing, forced by axioms C1 and C2 alone, implies that every C1-admissible observable on $\mathrm{Coh}(\Kfour)$ is invariant under $\Vfour$. Consequently the effective symmetry group of $\Kfour$ is
\[
G_{\mathrm{eff}} = \Sfour/\Vfour \;\cong\; \Sthree.
\]
This is an unconditional theorem of C1 and C2, with no auxiliary hypothesis.
\end{pdlbox}

\begin{pdlbox}
\textbf{Theorem (D58, Lemma L3).} The quotient $\Afour/\Vfour$ is cyclic of order 3:
\[
\Afour/\Vfour \;\cong\; \Zthree,
\]
and under the homomorphism $\rho \colon \Sfour \to \Sthree$ of the previous theorem, the image of $\Afour$ is exactly the alternating subgroup $\mathrm{Alt}(3) \subset \Sthree$ --- the centre-generating cyclic subgroup. D58 further proves, via the Cartan--Killing--Weyl correspondence and the $A_2$ root system realised in the rank-2 reduction forced by $\Re = 6$, that this $\Zthree$ is exactly the centre $Z(\SUthree)$ of the colour gauge group. The derivation of $\SUthree$ as an algebraic theorem of C1--C4 is complete.
\end{pdlbox}

\begin{pdlbox}
\textbf{Theorem (D59).} The fundamental representation $\mathbf{3}$ of $\SUthree$ is realised on the trace-zero plane $W = \{(a,b,c) \in \mathbb{C}^3 : a+b+c = 0\}$. The three axes of $W$ are canonically labelled, in an $\Sfour$-equivariant manner, by the three non-identity elements of $\Vfour$. The same three axes admit a second, dynamically motivated labelling by the triplet
\[
(T_2) = \{\kone, \ktwo, \kthree\} = \{9, 19, 168\},
\]
the indices of the three independent leakage cycles of D51--D52 associated with the cosmological constant $\Lambda_{\PDL}$. The matter orientation $\mathbf{3}$ (as opposed to $\bar{\mathbf{3}}$) is forced, without additional input, by the isospin asymmetry $\Dn = \ndd - \nuu = 4 > 0$, itself an unconditional theorem of C4.
\end{pdlbox}

\begin{remark}[What is and is not established]
The chain above establishes, as unconditional theorems of C1--C4, that: (i) the invariant~3 of N01 originates algebraically in the quotient $\Afour/\Vfour \cong \Zthree$; (ii) this $\Zthree$ is simultaneously the centre of the colour gauge group $\SUthree$ and the index set of the fundamental representation's three axes; (iii) the same three axes carry a second labelling by the three leakage-cycle indices $(\kone,\ktwo,\kthree)$ of the cosmological sector. What is \emph{not} established by this chain, and remains the subject of the present note, is any connection between this $\Zthree$ structure and the three \emph{fermion generations} of either framework. The \PDL\ corpus itself states this explicitly as OP-D59-1 (quark mass spectrum and generation structure from C1--C4 is open) and OP-D61-2 (three fermion generations from C1--C4 is open).
\end{remark}

% ============================================================
\section{The PDL leakage cycles: dynamical content of \texorpdfstring{$(T_2)$}{(T2)}}
\label{sec:leakage}
% ============================================================

The triplet $(T_2) = \{\kone, \ktwo, \kthree\}$ is not merely a relabelling of $\Vfour \setminus \{e\}$; it carries independent dynamical content established in D51--D52, which is the natural point of contact with the \OFN\ generation structure.

\begin{pdlbox}
\textbf{Theorem (D51, Lemma 1).} The first Betti number of the 1-skeleton of $\Kfour$ is $\beta_1(\Kfour) = |E| - |V| + 1 = 6 - 4 + 1 = 3$: a purely topological fact, independent of sign assignment or dynamics. There are exactly three independent leakage cycles.
\end{pdlbox}

\begin{pdlbox}
\textbf{Theorem (D51, Lemma 2--3; D52).} Axioms C1 and C3 force the lengths of admissible leakage cycles to be prime. The three leakage cycles of $\Kfour$ have prime exponents
\[
(\kone, \ktwo, \kthree) \;\longmapsto\; (p_{\kone}, p_{\ktwo}, p_{\kthree}) = (23, 67, 997),
\]
and the corresponding leakage bases are uniquely identified (D52) as $(1-\kappa)$, $\rval/\Rtot$, and $(1-\eta_L)$ respectively, each itself an unconditional theorem of C1--C4. The resulting cosmological leakage constant $C$ reproduces $\Lambda_{\mathrm{obs}}$ to $0.17$ parts per million (Planck 2020).
\end{pdlbox}

\begin{remark}
It is this triplet of \emph{prime-indexed, dynamically active} leakage cycles --- not merely the abstract group $\Zthree$ --- that is the most natural \PDL\ object to compare against three \OFN\ fermion generations, since fermion generations in any framework are expected to carry distinguishing dynamical content (e.g.\ mass hierarchy), not merely a permutation symmetry.
\end{remark}

% ============================================================
\section{Comparison table: PDL result \texorpdfstring{$\vert$}{|} OFN result \texorpdfstring{$\vert$}{|} status of identification}
\label{sec:comparison}
% ============================================================

Following the working conditions agreed for this note (Section~\ref{sec:conditions}), every cross-framework correspondence is recorded in the table below rather than in a shared equation. The \PDL\ column is complete; the \OFN\ column and the status column are to be completed jointly.

\begin{table}[H]
\centering
\small
\caption{Comparison table for the $\Zthree$ / three-generations correspondence and the five Session~67 connections. Status labels follow the \PDL\ epistemic convention: \textcolor{proofblue}{mathematical identity}, \textcolor{hypothesisorange}{structural echo / candidate analogy}, \textcolor{openproblemred}{open problem}.}
\label{tab:comparison}
\renewcommand{\arraystretch}{1.4}
\begin{tabular}{p{4.0cm}p{4.0cm}p{4.6cm}}
\toprule
\textbf{PDL result} & \textbf{OFN result} & \textbf{Status} \\
\midrule
$\beta_1(\Kfour) = 3$ (topological, D51 Lemma~1) &
$b_1(\Omegatf) = 3$ AND $b_1(G_H) = 3$ (N01 + Bachani CWS code, Sec.~\ref{sec:vacuum_manifold}) &
\textcolor{proofblue}{Mathematical identity (N01): same invariant, independent derivations.} \\[0.5em]
$\Afour/\Vfour \cong \Zthree \subset \mathrm{Sym}(\Kfour)$ (D58 Lemma~L3) &
$b_1(G_H) = 3$ cycles; the tripartition $\{A,B,C\}$ of $K_{2,2,2} = L(K_4)$ is the set of three perfect matchings of $K_4$ (Sec.~\ref{sec:session67_summary}) &
\textcolor{proofblue}{Mathematical identity: both are the same three-element set (matchings of $K_4$ = $V_4\setminus\{e\}$).} \\[0.5em]
$\Zthree = Z(\SUthree)$, centre of colour gauge group (D58) &
$13 = 8 \oplus 3 \oplus 1 \oplus 1$ (CWS code decomp.) encodes gauge group; $\Zthree$ is the centre of $\mathrm{SU}(3)_c$ &
\textcolor{hypothesisorange}{Shares integer 3; no structural identification claimed beyond this. OP-OFN-3 open.} \\[0.5em]
Three axes of $W$, labelled by $\Vfour\setminus\{e\}$ and by $(T_2)=\{\kone,\ktwo,\kthree\}$ (D59) &
Three cycles of $G_H$ with $b_1(G_H) = 3$; $S_3$ broken to $Z_2$ by holonomy distinction $\Phi_1=\Phi_2=2\pi/12$, $\Phi_3=5\pi/12$ &
\textcolor{hypothesisorange}{OP-OFN-1 partial resolution: two classes distinguished by holonomy; full three-generation separation open (OP-N02-2).} \\[0.5em]
$\phi$ in $\kappa = 310\phi/11017$ (D44, unconditional theorem) &
$\gamma_E = 3-\sqrt{5} = 4-2\phi$, spectral gap of $G_E$ (dodecahedron on $\Omega_{21}$, Theorem~\ref{thm:spectral_gap_GE}) &
\textcolor{proofblue}{Algebraic identity: $\phi + \gamma_E/2 = 2$ exact in $\mathbb{Q}(\sqrt{5})$ (Theorem~\ref{thm:algebraic_identity}). Note: $G_H$ has a different gap $\lambda_1\approx0.0804$, algebraically unrelated to $\phi$.} \\[0.5em]
$4$ vertices $+ 6$ edges of $K_4 = 10$ graph elements (D16a) &
$\tau = 90 = b_1^2 \times \dim P(1,3) = 9\times 10$ (Corollary~\ref{cor:structural_echo}); $\dim P(1,3)=4+6=10$ &
\textcolor{hypothesisorange}{Structural echo: OP-OFN-2 open.} \\[0.5em]
$(p_{k_1}, p_{k_2}, p_{k_3}) = (23,67,997)$, prime exponents of three leakage cycles (D51--D52) &
Three CP-pairs $+$ 13 unpaired states (Theorem~\ref{thm:cp_classification}); no OFN-side mass hierarchy from leakage primes yet identified &
\textcolor{hypothesisorange}{OP-OFN-1 and OP-N02-1 open.} \\[0.5em]
$n_u - 1 = p_{k_1} = 23$ (D47+D51; isolation 0.265\% / 378 pairs) &
No direct OFN counterpart identified at present &
\textcolor{hypothesisorange}{Speculative extension; OP-D63-3 open.} \\
\bottomrule
\end{tabular}
\end{table}

% ============================================================
\section{The OFN side: Spectral Structure of the Vacuum Manifold \texorpdfstring{$\Omega_{21}$}{Omega21}}
\label{sec:ssr}
% ============================================================

This section presents the \OFN-side results that complement the \PDL-side of this bridge note. All results concern the vacuum manifold $\Omega_{21}$, a 21-vertex subset of the 6-dimensional Boolean hypercube $Q_6$, and its spectral decomposition under the graph Laplacian.

\subsection{Definition of the Vacuum Manifold}
\label{sec:vacuum_manifold}

The vacuum manifold is defined as the unique 21-state subset of $Q_6$ selected by three cascading self-consistency conditions \cite{Bachani2026}:
\begin{enumerate}
    \item \textbf{Causal regularity:} vertex-transitive cubic subgraphs of $Q_6$ at Hamming distance 2.
    \item \textbf{Cage classification:} retention of graphs with girth $\geq 5$ (by Tutte's cage theorem, only the Petersen graph $P(10)$ and the dodecahedron $D(20)$ qualify).
    \item \textbf{Binary embedding:} self-consistent embedding $\Phi^*\colon \text{graph} \to Q_6$ preserving adjacency at Hamming distance 2. Only $D(20)$ admits such an embedding.
\end{enumerate}

The Hamming-isolated state $\sigma = 21 = (010101)_2$ completes the vacuum:
\[
\Omega_{21} = \Phi^*(D_{20}) \cup \{\sigma\} = \{0,1,3,4,7,8,9,12,15,16,19,21,27,31,35,42,43,48,52,56,63\}
\]

\subsection{Two Distinct Graphs on \texorpdfstring{$\Omega_{21}$}{Omega21}}
\label{sec:two_graphs}

\textbf{Critical clarification.} The vacuum manifold $\Omega_{21}$ carries \textit{two distinct graphs}, each with its own spectral gap. This distinction was not explicit in earlier \OFN\ literature and is essential for the \PDL--\OFN\ bridge.

\begin{definition}[Graph $G_E$ --- Dodecahedron]
\label{def:GE}
The graph $G_E$ is the dodecahedron $D(20)$, with vertex set $M_{20} = \Omega_{21} \setminus \{\sigma\}$ (the 20-vertex connected component) and edges defined by the graph structure of $D(20)$ embedded in $Q_6$ via the map $\Phi^*$ \cite{Bachani2026}. The graph $G_E$ is 3-regular with 30 edges. The state $\sigma = 21$ is \textit{not} part of $G_E$; it is an isolated vertex (degree 0) in the full graph on $\Omega_{21}$.
\end{definition}

\begin{definition}[Graph $G_H$ --- Induced Subgraph]
\label{def:GH}
The graph $G_H$ has both endpoints in $M_{20} = \Omega_{21} \setminus \{\sigma\}$, with edges at Hamming distance 1. It has 22 edges and is irregular, with vertex degrees $\{1, 2, 3, 4\}$.
\end{definition}

\textbf{Physical interpretation.} The two graphs correspond to different physical roles \cite{Bachani2026}: $G_E$ governs kinetic energy and propagation (dodecahedral topology); $G_H$ encodes internal topology and the three Betti cycles.

\subsection{Spectral Decomposition of \texorpdfstring{$G_H$}{GH}}
\label{sec:spectral_GH}

\begin{theorem}[Spectral Factorization over $\mathbb{Q}$]
\label{thm:spectral_factorization}
The characteristic polynomial of the Laplacian of the 20-vertex connected component of $G_H$ factorizes exactly over the rationals as:
\[
P_{20}(x) = x \cdot (x - 1)^2 \cdot (x - 3) \cdot P_{16}(x)
\]
where $P_{16}(x)$ is an irreducible polynomial of degree 16 over $\mathbb{Q}$ with constant term 600.
\end{theorem}

\begin{proof}
The factorization was verified computationally by exact integer arithmetic (Python, NumPy, NetworkX) and confirmed independently by both authors (July 2026). The irreducibility of $P_{16}(x)$ was verified by checking that it has no rational roots and does not factor over $\mathbb{Q}$ into polynomials of lower degree.
\end{proof}

\begin{corollary}[Exact Integer Eigenvalues]
\label{cor:integer_eigenvalues}
The Laplacian of $G_H$ has exact integer eigenvalues: $\lambda = 0$ (multiplicity 2, one per connected component), $\lambda = 1$ (multiplicity 2), and $\lambda = 3$ (multiplicity 1). The remaining 16 eigenvalues are algebraic numbers of degree 16, roots of $P_{16}(x)$.
\end{corollary}

\begin{theorem}[Spectral Gap of $G_H$]
\label{thm:spectral_gap_GH}
The spectral gap of $G_H$ (smallest nonzero eigenvalue of the Laplacian) is:
\[
\lambda_1 \approx 0.0804170036
\]
This is the smallest root of $P_{16}(x)$ and is an algebraic number of degree 16.
\end{theorem}

This value was computed independently by both authors, agreeing to 10 significant figures.

\subsection{Spectral Gap of \texorpdfstring{$G_E$}{GE} and the Algebraic Identity}
\label{sec:spectral_GE}

\begin{theorem}[Spectral Gap of $G_E$]
\label{thm:spectral_gap_GE}
The spectral gap of $G_E$ (the dodecahedron) is $\gamma_E = 3 - \sqrt{5} \approx 0.7639$, a known mathematical fact for the dodecahedron graph.
\end{theorem}

\begin{theorem}[Algebraic Identity in $\mathbb{Q}(\sqrt{5})$]
\label{thm:algebraic_identity}
The following identity holds exactly in $\mathbb{Q}(\sqrt{5})$:
\[
\phi + \frac{\gamma_E}{2} = 2
\]
where $\phi = (1 + \sqrt{5})/2$ is the golden ratio and $\gamma_E = 3 - \sqrt{5}$ is the spectral gap of $G_E$.
\end{theorem}

\begin{proof}
Direct computation: $\displaystyle\phi + \frac{\gamma_E}{2} = \frac{1 + \sqrt{5}}{2} + \frac{3 - \sqrt{5}}{2} = \frac{4}{2} = 2$.
\end{proof}

\begin{corollary}[Non-Identity for $G_H$]
\label{cor:non_identity_GH}
The analogous identity does \textit{not} hold for $G_H$: $\phi + \lambda_1/2 \approx 1.658 \neq 2$. The algebraic identity $\phi + \gamma/2 = 2$ is specific to the dodecahedron $G_E$, not the induced subgraph $G_H$.
\end{corollary}

\subsection{The Spanning Tree Invariant \texorpdfstring{$\tau = 90$}{tau = 90}}
\label{sec:spanning_trees}

\begin{theorem}[Spanning Trees of $G_H$]
\label{thm:spanning_trees}
The number of spanning trees of the 20-vertex connected component of $G_H$ is $\tau = 90$.
\end{theorem}

\begin{proof}
By the Matrix Tree Theorem (Kirchhoff), $\tau = \frac{1}{n}\prod_{i=1}^{n-1}\lambda_i$ where $n = 20$. From Theorem~\ref{thm:spectral_factorization}: $\prod_{i=1}^{19}\lambda_i = 1^2 \cdot 3 \cdot 600 = 1800$, so $\tau = 1800/20 = 90$.
\end{proof}

This was verified numerically to machine precision (product of nonzero eigenvalues = 1800.0000) by both authors.

\begin{corollary}[Structural Echo]
\label{cor:structural_echo}
The invariant $\tau = 90$ admits the factorization:
\[
\tau = 90 = 9 \times 10 = b_1^2 \times \dim P(1,3)
\]
where $b_1 = 3$ is the first Betti number of $G_H$ and $\dim P(1,3) = 10$ is the dimension of the Poincaré group (4 translations + 6 Lorentz generators). \textbf{Epistemic status:} this is a \textit{structural echo}, not a proven identity (see OP-OFN-2 below).
\end{corollary}

\subsection{CP-Pair Classification}
\label{sec:cp_pairs}

\begin{theorem}[CP-Pair Classification]
\label{thm:cp_classification}
Under the CP involution $s \mapsto 63 - s$, the 21 states of $\Omega_{21}$ decompose into exactly 4 CP-pairs (8 states) and 13 unpaired states. The 4 CP-pairs are: $(0,63)$, $(7,56)$, $(15,48)$, $(21,42)$.
\end{theorem}

\begin{proof}
The CP involution has no fixed points in $\mathbb{Z}$ (since $s = 63-s$ implies $s = 31.5$). Enumeration of all pairs $(s, 63-s)$ with both elements in $\Omega_{21}$ yields exactly the four pairs listed. The remaining 13 states have their CP-partner outside $\Omega_{21}$. Consistency: $4 \times 2 + 13 = 21$.
\end{proof}

\begin{remark}
The state $\sigma = 21 = (010101)_2$ is Hamming-isolated in $\Omega_{21}$ (degree 0 in $G_H$), but is \textit{not} a fixed point of the CP involution: its CP-partner is $42 = (101010)_2 \in \Omega_{21}$, forming the fourth CP-pair. This classification agrees with Bachani~\cite{Bachani2026}, who identifies the 8 states of the 4 CP-pairs as ``matter states'' and the 13 unpaired states as ``gauge channels'' --- where ``self-conjugate'' in Bachani's terminology refers to membership in a CP-pair within $\Omega_{21}$, not to fixed points of the involution.
\end{remark}

\subsection{Summary of the Five Session 67 Connections}
\label{sec:session67_summary}

The following five connections between \PDL\ and \OFN\ were established in July 2026:

\begin{enumerate}
    \item \textbf{Mathematical identity (tripartition):} The tripartition $\{A, B, C\}$ of $K_{2,2,2} = L(K_4)$ is exactly the set of three perfect matchings of $K_4$ (PDL D58 Lemma L2, D61). \textit{Status: mathematical identity.}
    \item \textbf{Algebraic identity in $\mathbb{Q}(\sqrt{5})$:} $\phi + \gamma_E/2 = 2$ (Theorem~\ref{thm:algebraic_identity}). \textit{Status: algebraic identity.}
    \item \textbf{Corrected automorphism group:} $\mathrm{Aut}(K_{2,2,2}) = S_2 \wr S_3$, order 48 ($O_h$), not $S_4$; Whitney's theorem does not apply to $K_4$. \textit{Status: corrected computation.}
    \item \textbf{Dimensional echo:} $K_4$ has 4 vertices and 6 edges, $4 + 6 = 10 = \dim P(1,3)$. \textit{Status: structural echo.}
    \item \textbf{Hadronic-cosmological identity:} $n_u - 1 = p_{k_1} = 23$ (D47 + D51), verified 13/13 at isolation 0.265\% against 378 pre-registered pairs. \textit{Status: verified identity.}
\end{enumerate}

\subsection{Epistemic Summary and Verification Record}
\label{sec:ofn_epistemic}

\textbf{Unconditional theorems:} Spectral factorization $P_{20}(x)$ (Theorem~\ref{thm:spectral_factorization}); spectral gap $\lambda_1 \approx 0.0804$ of $G_H$ (Theorem~\ref{thm:spectral_gap_GH}); spectral gap $\gamma_E = 3 - \sqrt{5}$ of $G_E$ (Theorem~\ref{thm:spectral_gap_GE}); algebraic identity $\phi + \gamma_E/2 = 2$ (Theorem~\ref{thm:algebraic_identity}); spanning tree invariant $\tau = 90$ (Theorem~\ref{thm:spanning_trees}); CP-pair classification (Theorem~\ref{thm:cp_classification}).

\textbf{Structural echoes (not proven identities):} $\tau = 90 = b_1^2 \times \dim P(1,3)$; $4 + 6 = 10 = \dim P(1,3)$.

All computations were performed by exact integer arithmetic (Python, NumPy, NetworkX) and independently verified by both authors in July 2026. Scripts are available upon request.

% ============================================================

\section{Speculative extension: the identity \texorpdfstring{$\nuu - 1 = \pkone = 23$}{nu - 1 = p k1 = 23}}
\label{sec:speculative}
% ============================================================

\begin{speculativebox}
This section is presented at a lower epistemic register than Sections~\ref{sec:pdl_side}--\ref{sec:leakage}. The chain $\Afour/\Vfour \cong \Zthree \to (T_2) \to$ axes of $W$ is a sequence of unconditional theorems. The material below is a numerically verified \emph{coincidence} whose structural origin is an open problem (OP-D63-3) of the \PDL\ corpus. It is included here because it independently connects the same leakage index $\kone$ to a second \PDL\ sector (hadronic masses), and may therefore be of interest as an additional, lower-confidence data point for the \OFN\ comparison --- not as a second derivation chain of comparable strength to Sections~\ref{sec:pdl_side}--\ref{sec:leakage}.
\end{speculativebox}

\subsection{The identity}

D63~\cite{D63} records the following exact integer identity between two quantities derived from independent \PDL\ mechanisms:
\begin{equation}
\nuu - 1 \;=\; \pkone \;=\; 23,
\label{eq:identity}
\end{equation}
where $\nuu = 24$ is the up-quark valence multiplicity of the proton quintuplet (D47, forced by the quasi-completeness constraint $3\nuu^2 + 5\nuu - 1848 = 0$, discriminant $149^2$), and $\pkone = 23$ is the first leakage-cycle prime of the cosmological sector (D51, with $\kone = \ndd - \nuu + \Re - 1 = 9$ and $\pkone$ the $9^{\mathrm{th}}$ prime).

\subsection{Verification and isolation test}

The identity~\eqref{eq:identity} and its statistical isolation within the natural reference class of simple linear expressions in the quintuplet integers were verified independently in the script \\ \texttt{PDL\_N02\_identity\_lockdown\_v2\_reinforced.py}~\cite{PDLN02lockdown}. The script: (i) recomputes $\nuu = 24$ from the quasi-completeness constraint independently of its statement in D63; (ii) recomputes $\kone = 9$ and $\pkone = 23$ from the leakage-index definition of D51; (iii) verifies the identity~\eqref{eq:identity} by exact integer arithmetic; (iv) tests a pre-registered family of 378 candidate (index, value) expression pairs built from simple linear combinations of $(\nuu, \ndd, \Re, \Dn)$, finding that only $1$ pair --- the documented identity itself --- reproduces an exact prime-index match (isolation fraction $0.265\%$); (v) cross-checks this rate against a randomised null model over the same coefficient ranges ($5000$ draws, match rate $0.059\%$), confirming the pre-registered family was not tuned to produce an artificially low rate. All checks passed ($13/13$).

\begin{remark}[Epistemic status]
This result supports treating the identity~\eqref{eq:identity} as a non-trivial structural coincidence --- it is not a generic feature of prime density, nor an artefact of a cherry-picked comparison family --- but it does \emph{not} constitute a derivation from C1--C4. The structural origin of the identity remains OP-D63-3, open in the \PDL\ corpus. No claim is made here regarding any \OFN-side counterpart of this specific identity; this section is offered to Oleg for consideration only.
\end{remark}

% ============================================================
\section{Working conditions for this note}
\label{sec:conditions}
% ============================================================

Consistent with the conditions agreed by both authors prior to drafting, and modelled on the epistemic discipline of N01:

\begin{enumerate}[leftmargin=2em]
\item Every correspondence between a \PDL\ object and an \OFN\ object is recorded as a \emph{candidate structural analogy} (Table~\ref{tab:comparison}) until a derivation common to both frameworks exists.
\item The structural self-reference operator $\Ssr$ is introduced strictly as an \OFN-side concept originating in~\cite{DTOCOFN}, clearly attributed, and is not identified with any \PDL\ axiom or theorem.
\item No equation in this note mixes objects native to both frameworks; cross-framework statements are confined to Table~\ref{tab:comparison} and to remarks explicitly labelled as candidate analogies.
\item Results are stratified by epistemic status throughout: \textcolor{proofblue}{unconditional theorem}, \textcolor{conjecturegreen}{corroborated conjecture}, \textcolor{hypothesisorange}{candidate structural analogy}, or open problem, following the convention of the \PDL\ corpus~\cite{DMv29}.
\end{enumerate}

% ============================================================
\section{Open problems}
\label{sec:open}
% ============================================================

\begin{openproblem}[OP-OFN-1 --- restated]
Formally identify the $\Zthree$ structure of Section~\ref{sec:pdl_side} with the three fermion generations of \OFN. The holonomy mechanism of Section~\ref{sec:ssr} (brisure $S_3\to Z_2$, $\Phi_1=\Phi_2\neq\Phi_3$) provides a partial resolution: two classes are distinguished dynamically. The full separation of all three generations and the mass hierarchy remain open.
\end{openproblem}

\begin{openproblem}[OP-D63-3 --- restated]
Determine whether the identity $\nuu - 1 = \pkone = 23$ (Section~\ref{sec:speculative}) is forced by C1--C4 or is a coincidence at the current level of approximation. \textbf{Entry point:} D47, D51, D63, present note.
\end{openproblem}

\begin{openproblem}[New --- OP-N02-1]
Determine whether the dynamical content of the three leakage cycles (prime exponents $23, 67, 997$, Section~\ref{sec:leakage}) admits any \OFN-side counterpart distinguishing the three fermion generations dynamically (e.g.\ a mass-hierarchy mechanism), beyond the permutation symmetry $\Zthree$ common to both frameworks. \textbf{Entry point:} D51, D52, $G_H$ holonomy (Section~\ref{sec:ssr}).
\end{openproblem}

\begin{openproblem}[New --- OP-N02-2]
Break the residual $Z_2$ symmetry between cycles $C_1$ and $C_2$ of $G_H$ (which share equal holonomies $\Phi_1 = \Phi_2 = 2\pi/12$, a topological consequence of their identical edge-class composition). This requires a more refined phase rule $\Theta$ accounting for additional structure (e.g.\ Hamming weights of specific vertices or higher-order holonomy corrections). \textbf{Entry point:} Section~\ref{sec:ssr}, $G_H(\Omega_{21})$.
\end{openproblem}

\begin{openproblem}[New --- OP-OFN-2]
Determine whether the factorization $\tau = 90 = b_1^2 \times \dim P(1,3)$ reflects a deep structural connection between the spanning tree structure of $G_H$ and the Poincaré group, or is a coincidence of small numbers. \textbf{Entry point:} Corollary~\ref{cor:structural_echo}, Section~\ref{sec:ssr}.
\end{openproblem}

\begin{openproblem}[New --- OP-OFN-3]
Derive $\mathrm{SU}(3)\times\mathrm{SU}(2)\times\mathrm{U}(1)$ from the \OFN\ decomposition $8\oplus 3\oplus 1\oplus 1$ of the CP-crossing channels (Theorem~\ref{thm:cp_classification}). The \PDL\ side has achieved this derivation (D46, D60, D58, D59); the \OFN-side derivation remains open.
\end{openproblem}

\begin{openproblem}[New --- OP-OFN-4]
Derive the topological impedance $\alpha$ in the redshift formula $(1+z) = \exp(\alpha d)$ from the spectral structure of $\Omega_{21}$. Candidate: $\alpha = f(\lambda_1, \tau)$ where $\lambda_1\approx0.0804$ and $\tau=90$, but $f$ must also encode the hierarchy between the fundamental network scale and the cosmological horizon ($\alpha_{\mathrm{phys}}\approx H_0\cdot t_D\sim10^{-61}$, requiring additional structure). \textbf{Epistemic status:} conjecture, not theorem.
\end{openproblem}

% ============================================================
\section*{Acknowledgements}
% ============================================================
The authors thank each other for the ongoing collaboration initiated in N01. [Further acknowledgements to be added jointly.]

\clearpage

% ============================================================
\bibliographystyle{unsrt}
\bibliography{B2_references}
% ============================================================

\end{document}
